% !TEX root = ../main.tex
\chapter{Reverse-Mode Automatic Differentiation Theory}
\label{ch:ad_theory}

Every framework that trains neural networks computes gradients.
The how-of-computing varies more than one might think: symbolic
differentiation (Mathematica-style), numerical finite differences,
forward-mode automatic differentiation, and reverse-mode automatic
differentiation are four distinct techniques with different
complexity profiles and different failure modes
\cite{baydin_automatic_2018,griewank_evaluating_2008}. \Lucid, like every
other production ML framework, uses reverse-mode automatic
differentiation as its core gradient algorithm. This chapter
explains why, what reverse-mode AD is, what guarantees it provides,
and what its costs are.

We begin with the theory and end with the practical consequences
that the implementation chapters
(\chref{ch:autograd_engine}, \chref{ch:custom_functions},
\chref{ch:gradient_checkpointing}, \chref{ch:higher_order})
elaborate.

\section{Four Ways to Compute a Derivative}
\label{sec:ad_theory:four_ways}

Consider a scalar function $L(x_1, x_2, \ldots, x_N)$ ---
``$L$'' for loss --- and the goal of computing
$\partial L / \partial x_i$ for every input. There are four
approaches.

\subsection{Symbolic differentiation}
\label{ssec:ad_theory:symbolic}

Symbolic differentiation manipulates the algebraic expression for
$L$ to derive a new expression for $\partial L / \partial x_i$.
Mathematica and SymPy do this. The output is exact and
human-readable.

The pathology: expression swell. For a deep composition like
$L = \sigma(W_n \sigma(W_{n-1} \sigma(\cdots W_1 x)))$, the
symbolic derivative has $O(2^n)$ terms even after simplification.
Symbolic differentiation is unsuitable for any function with more
than a few dozen primitive ops, which means unsuitable for
neural networks.

\subsection{Numerical differentiation}
\label{ssec:ad_theory:numerical}

Numerical differentiation approximates the derivative by finite
differences:
\[
\frac{\partial L}{\partial x_i} \approx
\frac{L(x_1, \ldots, x_i + h, \ldots, x_N) - L(x_1, \ldots, x_i,
\ldots, x_N)}{h}.
\]
For $N$ inputs, this requires $N + 1$ evaluations of $L$ (or $2N$
for the more accurate two-sided variant). The accuracy is bounded
by $O(h)$ truncation error and $O(\epsilon / h)$ round-off error
(where $\epsilon$ is machine precision), giving a best case of
$O(\sqrt{\epsilon}) \approx 10^{-4}$ for FP32.

The cost is prohibitive for ML: a 100M-parameter model requires
$10^8$ evaluations of the forward pass per gradient computation.
Finite differences are useful only as a correctness check
(gradcheck; \chref{ch:gradcheck}), not as a training-time
algorithm.

\subsection{Forward-mode AD}
\label{ssec:ad_theory:forward_mode}

Forward-mode AD propagates derivatives alongside values through
the computation. For each input $x_i$, we maintain a pair $(v,
\dot{v})$ where $v$ is the value and $\dot{v}$ is the derivative
of $v$ with respect to $x_i$. The chain rule says that for any
function $y = f(v)$,
\[
\dot{y} = f'(v) \cdot \dot{v}.
\]

For $N$ inputs and one output, forward mode requires $N$ passes
to compute the full gradient. Each pass costs roughly the same as
the original forward computation, so the total cost is $N$ times
forward. Like finite differences, this is prohibitive for ML.

Forward mode is the right choice when the number of \emph{outputs}
greatly exceeds the number of inputs (e.g., a Jacobian of a $1
\to 1000$ function). For ML, where the loss is a scalar and the
parameters are millions, the opposite is true.

\subsection{Reverse-mode AD}
\label{ssec:ad_theory:reverse_mode}

Reverse-mode AD propagates derivatives \emph{backward} from the
output. For each intermediate $v$ in the computation, we maintain
a \emph{cotangent}
\[
\bar{v} \;=\; \frac{\partial L}{\partial v}.
\]
The chain rule, applied in reverse, says that if $y = f(v)$
contributed to $L$, then $\bar{v}$ accumulates a term
\[
\bar{v} \;\mathrel{+}=\; \left(\frac{\partial f}{\partial
v}\right)^{\!\top} \bar{y}.
\]
For a multivariate $f : \mathbb{R}^n \to \mathbb{R}^m$ with
Jacobian $J = \partial f / \partial v$, the contribution is the
vector-Jacobian product $J^\top \bar{y}$. The transpose is
essential: the Jacobian itself ($J$) maps tangents forward; its
transpose maps cotangents backward.

For a single scalar output and $N$ inputs, reverse mode computes
the full gradient in a single backward pass, costing roughly 2--3
times the forward pass. This is the algorithmic basis of every
modern ML framework's training loop.

The trade-off: reverse mode requires storing intermediate values
(activations) from the forward pass, because the backward pass
needs them. Memory grows with the depth and width of the forward
computation. This is why training a deep network requires
substantially more memory than inference, and why gradient
checkpointing (\chref{ch:gradient_checkpointing}) exists.

\tabref{tab:ad_theory:comparison} summarises the four.

\begin{table}[t]
\centering
\small
\begin{tabular}{llll}
\toprule
Method & Cost (per gradient) & Accuracy & Best when \\
\midrule
Symbolic     & $O(\text{expr swell})$ & exact          & expression is short \\
Numerical    & $O(N \cdot T_{\text{fwd}})$ & $\sqrt{\epsilon}$  & correctness check \\
Forward AD   & $O(N \cdot T_{\text{fwd}})$ & exact          & many outputs, few inputs \\
Reverse AD   & $O(T_{\text{fwd}})$         & exact          & many inputs, few outputs (ML) \\
\bottomrule
\end{tabular}
\caption[AD method comparison.]{Cost and accuracy trade-offs for
the four differentiation techniques. For ML, where a model has
millions of parameters and a single scalar loss, reverse-mode AD
is the only viable option.}
\label{tab:ad_theory:comparison}
\end{table}

\section{The Chain Rule in Reverse}
\label{sec:ad_theory:chain_rule_reverse}

The mathematical basis is the multivariate chain rule, applied to
a directed acyclic graph of operations.

\subsection{Single-variable chain rule}
\label{ssec:ad_theory:chain_single}

For a composition $L = h(g(f(x)))$, the derivative is
\[
\frac{\partial L}{\partial x} = h'(g(f(x))) \cdot g'(f(x)) \cdot
f'(x).
\]

Reading right-to-left is forward mode: start with $\dot{x} = 1$,
propagate through $f$, $g$, $h$. Reading left-to-right is reverse
mode: start with $\bar{L} = 1$, propagate through $h$, $g$, $f$.

\subsection{Multi-variable chain rule}
\label{ssec:ad_theory:chain_multi}

For a multivariate function $y = f(v_1, v_2, \ldots, v_k)$, the
chain rule contribution to the cotangent of each input is
\[
\bar{v_i} \;\mathrel{+}=\; \left(\frac{\partial f}{\partial
v_i}\right)^{\!\top} \bar{y}.
\]

If $v_i$ feeds into multiple downstream operations
$y_1 = f_1(v_i, \cdot), y_2 = f_2(v_i, \cdot), \ldots, y_m =
f_m(v_i, \cdot)$, every path contributes, so
\[
\bar{v_i} \;=\; \sum_{j=1}^{m} \left(\frac{\partial f_j}{\partial
v_i}\right)^{\!\top} \bar{y_j}.
\]
The $+=$ is the \emph{accumulation}: each path's contribution
adds to $\bar{v_i}$. This is the mechanism by which a parameter
used in multiple places (the shared weight matrix of a recurrent
network, the embedding shared between encoder and decoder) gets
correct gradient summed across all its uses.

\subsection{Computation graph as DAG}
\label{ssec:ad_theory:dag}

The computation is naturally a DAG: nodes are operations, edges
are tensor flow. Forward execution walks the DAG topologically
from inputs to outputs. Backward execution walks the DAG in
reverse topological order from the loss to the parameters.

\figref{fig:ad_theory:dag} shows a small example.

\begin{figure}[t]
\centering
\begin{adjustbox}{max width=\textwidth, center}
\begin{tikzpicture}[
  inp/.style={draw, circle, fill=lucid.note.bg, minimum size=0.9cm,
              font=\small\sffamily},
  op/.style={draw, rounded corners=2pt, fill=lucid.info.bg,
             minimum width=1.4cm, minimum height=0.8cm,
             font=\small\sffamily, align=center},
  out/.style={draw, circle, fill=lucid.ok.bg, minimum size=0.9cm,
              font=\small\sffamily\bfseries},
  fwd/.style={->, >=stealth, thick, color=lucid.accent.dark,
              shorten >=2pt, shorten <=2pt},
  lbl/.style={font=\footnotesize\itshape, color=lucid.muted,
              fill=white, inner sep=2pt},
]
  \node[inp] (x1) at (0, 2)    {$x_1$};
  \node[inp] (x2) at (0, 0)    {$x_2$};
  \node[op]  (m)  at (3, 1)    {$\times$};
  \node[op]  (a)  at (6, 0.5)  {$+$};
  \node[op]  (s)  at (9, 0.5)  {$\sigma$};
  \node[out] (L)  at (12, 0.5) {$L$};

  \draw[fwd] (x1) -- (m);
  \draw[fwd] (x2) -- (m);
  \draw[fwd] (x2) to[bend right=20] (a.south west);
  \draw[fwd] (m) -- (a);
  \draw[fwd] (a) -- (s);
  \draw[fwd] (s) -- (L);

  \node[lbl] at (12, -0.5) {$\bar{L} = 1$ (seed)};
  \node[lbl, align=left, anchor=west] at (0, -1.3)
    {$\bar{x}_2 \mathrel{+}=$ contribution from $\times$\\
     $\bar{x}_2 \mathrel{+}=$ contribution from $+$};
\end{tikzpicture}
\end{adjustbox}
\caption[Forward and backward through a DAG.]{A small
computation: $L = \sigma(x_1 x_2 + x_2)$. Solid arrows show
forward execution (inputs $\to$ output). The backward pass
traverses the same DAG in reverse, accumulating cotangents.
$x_2$ contributes through two paths (the multiplication and the
addition), so its cotangent accumulates from both contributions
--- the multi-variable chain rule in action.}
\label{fig:ad_theory:dag}
\end{figure}

\section{Vector-Jacobian Products}
\label{sec:ad_theory:vjp}

Reverse-mode AD is sometimes described in terms of
\emph{vector-Jacobian products} (VJPs). The reformulation is
useful because it makes the per-op backward function's signature
explicit and connects the theory to the implementation
(\chref{ch:custom_functions}).

\subsection{The Jacobian}
\label{ssec:ad_theory:jacobian}

For a function $f : \mathbb{R}^n \to \mathbb{R}^m$, the Jacobian
$J$ is the $m \times n$ matrix
\[
J_{ij} = \frac{\partial f_i}{\partial x_j}.
\]
The Jacobian is the linear approximation of $f$ near a point.

For a scalar-output loss $L = g(f(x))$, the chain rule says
\[
\nabla_x L = J^\top \cdot \nabla_{f(x)} g.
\]

The right-hand side is a Jacobian \emph{transpose} times a
vector. Computing the full $J$ explicitly is wasteful for the
common $m \ll n$ case; instead we compute the product $J^\top
\bar{y}$ directly. This is the \emph{vector-Jacobian product}
(VJP).

\subsection{VJP as the per-op backward}
\label{ssec:ad_theory:vjp_per_op}

Each primitive operation in \Lucid\ ships a VJP function that
takes the output cotangent and produces the input cotangents:
\[
\text{backward}(\bar{y};\ \text{saved}) = (\bar{x_1}, \bar{x_2},
\ldots).
\]

Concrete examples for the most common ops, derived from the
forward expression $y = f(\cdot)$ and the rule $\bar{x_i} =
(\partial f / \partial x_i)^\top \bar{y}$:

\begin{align*}
y = a + b              &: \quad (\bar{a}, \bar{b}) = (\bar{y},\, \bar{y}) \\
y = a \odot b          &: \quad (\bar{a}, \bar{b}) = (b \odot \bar{y},\, a \odot \bar{y}) \\
Y = A B                &: \quad (\bar{A}, \bar{B}) = (\bar{Y} B^\top,\, A^\top \bar{Y}) \\
y = \mathrm{ReLU}(x)   &: \quad \bar{x} = \mathbb{1}[x > 0] \odot \bar{y} \\
y = \sigma(x)          &: \quad \bar{x} = \sigma(x)(1-\sigma(x)) \odot \bar{y} \\
y = \tanh(x)           &: \quad \bar{x} = (1 - \tanh^2(x)) \odot \bar{y} \\
y = \mathrm{softmax}(x)&: \quad \bar{x} = y \odot (\bar{y} - (y \cdot \bar{y})) \\
y = \mathrm{logsumexp}(x) &: \quad \bar{x} = \mathrm{softmax}(x) \cdot \bar{y}
\end{align*}

The Custom Function API (\chref{ch:custom_functions}) lets the
user supply a VJP for an op that the framework does not know
about; the autograd engine threads it into the same chain as the
built-in VJPs.

The Custom Function API (\chref{ch:custom_functions}) lets the
user supply a VJP for an op that the framework does not know
about; the autograd engine threads it into the same chain as the
built-in VJPs.

\subsection{The Jacobian-vector product (JVP) is forward mode}
\label{ssec:ad_theory:jvp}

The mirror image is the Jacobian-vector product:
\[
\text{forward AD}(\dot{x}) = J \cdot \dot{x}.
\]
This is what forward mode computes per step. \Lucid\ exposes the
JVP through \code{lucid.func.jvp}
(\chref{ch:higher_order}); it is used in higher-order constructs
like Hessian-vector products.

\section{Reverse Mode's Complexity}
\label{sec:ad_theory:complexity}

The famous \emph{cheap gradient principle} of AD: a reverse-mode
gradient of any function costs at most a small constant ($\leq
5$, typically 2--3) times the forward function. The proof is by
induction on the computation graph: each primitive's VJP costs
$O(\text{primitive's forward})$, and the sum over the graph is
proportional to the forward sum.

\subsection{The 3x bound in practice}
\label{ssec:ad_theory:three_x}

For \Lucid's typical workloads, the measured backward-to-forward
ratio is in the range 2.0 to 3.5. The breakdown:

\begin{itemize}
  \item Each primitive's backward kernel has approximately the
    same flop count as its forward (1x).
  \item Saving activations for backward incurs no extra flops but
    does incur memory writes during forward and reads during
    backward, roughly equal to the forward's memory traffic
    (1x in memory bandwidth).
  \item The framework overhead of constructing the backward graph
    is on the order of a few microseconds per op.
\end{itemize}

The 3x heuristic explains why training is roughly 3$\times$
slower than inference for the same model on the same hardware.

\subsection{Memory cost}
\label{ssec:ad_theory:memory_cost}

The cheap gradient principle is about time, not memory. Memory
cost is the price reverse mode pays: every intermediate value
needed by a downstream backward is held alive from the forward
that produces it until the backward that consumes it.

For a deep network with $L$ layers, the activation memory is
$O(L)$ in the worst case. Gradient checkpointing
(\chref{ch:gradient_checkpointing}) trades: hold only
$O(\sqrt{L})$ activations, recompute the rest from checkpoints
during backward, paying an extra factor of $\sqrt{L}$ in forward
time.

\section{Practical Implications for Lucid}
\label{sec:ad_theory:implications}

The theory translates to a small set of implementation choices
that the rest of \partref{part:autograd} elaborates.

\subsection{Every primitive has a VJP}
\label{ssec:ad_theory:every_vjp}

Every differentiable op in \Lucid\ has a registered backward
function (\chref{ch:op_registry}'s \code{BackwardSpec}). The
autograd engine constructs a backward node from the schema at
forward-call time and walks the chain at \code{backward()} time.

For ops that the user defines, the Custom Function API
(\chref{ch:custom_functions}) is the mechanism for supplying a
VJP.

\subsection{Saved tensors are explicit}
\label{ssec:ad_theory:saved_explicit}

The backward function for a primitive needs some forward
quantities to compute its VJP. The op's schema declares which:
the input, the output, both, or neither. The autograd engine
saves only what is declared, which keeps memory usage bounded.

The user-side Custom Function API exposes the same mechanism
through \code{ctx.save\_for\_backward(...)}.

\subsection{Accumulation, not replacement}
\label{ssec:ad_theory:accumulation}

The chain rule says cotangents \emph{accumulate}. If a tensor $v$
is consumed by multiple downstream operations, each contributes
to $\bar{v}$. The autograd engine handles this by initialising
$\bar{v}$ to zero and adding each downstream contribution. The
\code{AccumulateGrad} node on leaf tensors performs the final
accumulation into \code{tensor.grad}.

A consequence: \code{tensor.grad} is not reset between training
steps automatically. The user must call
\code{optimizer.zero\_grad()} (or equivalent) at the start of
each step, or gradients from previous steps will continue to
accumulate.

\subsection{Determinism considerations}
\label{ssec:ad_theory:determinism}

Reverse mode accumulates cotangents in an order determined by the
DAG topology, which is deterministic. But the floating-point
addition is not associative, so a parallel implementation can
produce different bit-level results on different runs.

\Lucid's \code{set\_deterministic(True)} mode
(\chref{ch:reproducibility}) forces a serial accumulation order
for scatter-add and similar non-associative reductions. The cost
is performance; the benefit is bit-reproducible gradients.

\section{Higher-Order Differentiation}
\label{sec:ad_theory:higher_order}

\Lucid\ also supports higher-order derivatives --- the gradient of
a gradient, or the Hessian, or arbitrary compositions of VJP and
JVP. The mechanism is to make the autograd engine itself
differentiable: backward passes construct new graphs that can
themselves be differentiated.

\subsection{Differentiable backward}
\label{ssec:ad_theory:diff_back}

When \code{backward(create\_graph=True)} is called, the autograd
engine records the backward operations themselves into a new
graph. A subsequent \code{backward()} on the resulting gradient
computes the second-order derivative.

This is the mechanism that powers \code{lucid.func.hessian} and
related higher-order constructs.

\subsection{Composing transforms}
\label{ssec:ad_theory:composing}

The functional transforms (\code{vmap}, \code{grad}, \code{jvp},
\code{vjp}, \code{jacrev}, \code{jacfwd}) all build on the
underlying reverse-mode infrastructure. They compose: \code{grad
o vmap} computes per-sample gradients in parallel; \code{jacrev
o jacrev} computes the Hessian via two reverse passes;
\code{jacfwd o jacrev} computes the Hessian via mixed mode. The
chapter for these is \chref{ch:higher_order}.

\section{What Reverse Mode Does Not Do}
\label{sec:ad_theory:not_do}

A short list of misconceptions.

\subsection{It does not give symbolic expressions}
\label{ssec:ad_theory:no_symbolic}

The output of reverse mode is a numerical gradient at a specific
input point, not an algebraic expression. The output of $\nabla L$
at the point $x = (1.0, 2.0)$ is a vector of floats. \Lucid\ has
no way to give the user a symbolic expression for $\nabla L$ as a
function of $x$; the cost would be prohibitive.

\subsection{It does not give exact second-order information for free}
\label{ssec:ad_theory:no_free_second}

A single reverse pass gives first-order gradient. The Hessian
requires a second pass (or a different algorithm; see
\chref{ch:higher_order}). The Hessian of an $N$-parameter model
has $N^2$ entries and is rarely computed in full; the practical
substitutes are Hessian-vector products and quasi-Newton
approximations.

\subsection{It does not handle non-differentiable functions silently}
\label{ssec:ad_theory:no_nondiff}

Some operations are non-differentiable at certain points (ReLU at
0, abs at 0, argmax everywhere). Convention assigns these a
subgradient (typically zero), which \Lucid\ does. The result is a
valid subgradient but not the true gradient at the
non-differentiable point. Most training procedures are robust to
this; the user should be aware.

\section{Summary and Forward References}
\label{sec:ad_theory:summary}

Reverse-mode AD is the algorithm at the heart of \Lucid's
training loop. It computes a scalar-output gradient with respect
to any number of inputs in time proportional to a single forward
pass, at the cost of storing intermediate activations from the
forward. Every primitive in \Lucid\ ships a backward function
(its VJP); the engine assembles these into a per-call backward
graph; \code{backward()} walks the graph in reverse topological
order, accumulating cotangents.

The next chapter, \chref{ch:autograd_engine}, describes the
engine that implements this algorithm in \Lucid's C++ core.
\chref{ch:custom_functions} treats the user-facing extension
mechanism. \chref{ch:gradient_checkpointing} treats the
memory-vs-time trade-off. \chref{ch:higher_order} treats the
functional transforms that compose on top.
\chref{ch:gradcheck} treats finite-difference verification.

\begin{lucidnote}
For the user: \code{tensor.backward()} is the entry point. Behind
it sits the entire machinery this chapter described. The
machinery is invisible during a training step --- the user calls
\code{loss.backward()}, the gradients appear in the parameters'
\code{.grad} fields, and the optimiser takes a step. The
machinery becomes visible only when something goes wrong (a NaN
in a gradient, an unexpected zero, a missing backward for a
custom op).
\end{lucidnote}
